% <<<PDFPAGE 50>>>
\chapter{推理的未来}

\section{扩散式语言模型}

逐个词元地自回归生成文本，只是生成范式之一。扩散式语言模型（diffusion-based language models），如 MDLM、Plaid 和 LLaDA，把扩散过程应用于离散词元序列。前向过程逐渐遮盖词元，模型则学习去噪（denoising），也就是去除掩码以恢复词元。与自回归模型不同，扩散式大语言模型可以并行或半并行地生成所有词元，因此有望在生成过程中显著提高硬件利用率。

文本扩散模型在生成质量上还没有达到自回归模型的水平。迭代去噪的成本也可能比预期更高，通常需要 10–100 步，每一步都处理全部词元。不过，LLaDA 等近期模型已经明显缩小了质量差距。

\section{多词元预测}

Meta 的多词元预测（multi-token prediction，Gloeckle 等，2024）使用多个输出头来训练模型，让模型同时预测接下来的 $k$ 个词元，而非仅预测下一个词元。推理时，可以像 Medusa 的自推测解码（self-speculative decoding）那样，每一步生成多个词元，也可以用多个预测头为下游任务得到更好的内部表征。以四词元预测训练的模型比采用标准下一词元预测的模型具有更高的样本效率和更好的表征质量。

\section{神经形态计算与光子计算}

神经形态芯片（neuromorphic chips）如 Intel Loihi 2 和 IBM NorthPole，采用类似神经元脉冲的计算方式，对稀疏激活尤其节能，单次运算的能效比在 GPU 上运行深度学习高出几个数量级。IBM NorthPole 芯片在稠密神经网络推理中可达 22~TOPS/W，而 GPU 约为 2~TOPS/W。随着大语言模型的架构越来越多地利用稀疏性，神经形态硬件可能也会派上用场。

光子矩阵乘法器（photonic matrix multiplier），如 Lightmatter 和 Luminous Computing 的产品，利用光完成矩阵–向量乘法，理论上可获得无限的带宽，且每次乘法的能耗接近于零。光子计算尚处早期阶段，但如果加载权重几乎无需付出代价，它就有可能从根本上改变推理受访存带宽限制的局面。

\section{推理优先的设计理念}

这一领域越来越重视从模型设计之初就考虑推理需求。
\begin{itemize}
  \item \textbf{提高 KV 缓存效率的架构}。MQA、GQA、MLA 和 SSM 层都能减少 KV 缓存的占用。
  \item \textbf{激活稀疏性（activation sparsity）}。以 ReLU 或 top-$K$ 激活函数训练的模型会产生稀疏激活，可借此加快前馈网络的计算。
% <<<PDFPAGE 51>>>
  \item \textbf{混合深度（Mixture of Depths，MoD）}（Raposo 等，2024）。动态分配计算，词元按难度只在部分层参与注意力运算。
  \item \textbf{提前退出}。对易于处理的词元，根据置信度在中间层结束计算。
  \item \textbf{字节级模型（byte-level models）}。直接处理字节，免去子词切分的开销瓶颈，并改善多语种表现。
\end{itemize}

\begin{interview}
\question{1}{扩散式语言模型生成文本的方式与自回归模型有何区别？对推理有什么影响？}
自回归模型每次预测一个词元，并以此前的所有词元为条件。这一过程本质上是串行的，且受访存带宽限制。扩散式语言模型同时处理整个序列。它从一个达到目标长度、所有位置均被掩码遮盖的序列开始，在 $T$ 步内用模型预测哪些词元应解除掩码，逐步去噪。这样做对推理有以下影响。（1）\textbf{并行性}。每个去噪步同时处理所有位置，能比自回归解码更充分地利用张量核心。（2）\textbf{可控性}。扩散模型可以自然地填充遮盖的部分（用于补全或编辑），无须从头生成。（3）\textbf{成本结构}。对同样长度为 $N$ 的输出，扩散模型要进行 $T\times N$ 次词元预测，自回归模型要进行 $N$ 次，扩散模型的成本可能更高。（4）\textbf{长度灵活性}。要预先知道输出长度，或将其设为超参数。LLaDA 等模型使用 64–256 步，即 $T$ 取这个范围内的值。用蒸馏可以减至 4–8 步。

\question{2}{什么是混合深度？它如何按需分配推理计算？}
在混合深度框架中（Raposo 等，2024），词元可以跳过某些 Transformer 层。每层有一个路由器，决定词元是交给本层的注意力和前馈网络处理，还是沿残差连接直接传递，以跳过本层。简单、容易预测的词元可以跳过一些层，复杂、不确定的词元则接受完整处理。推理时的计算量因而可以随输入变化。一条简单的事实查询可能只有效使用了 30\% 的模型层，一条复杂的推理查询则用到所有层。提前退出会在某一层停止计算，MoD 则让跳过某层的词元继续经过后续层。训练时需要仔细设计路由，通常在各层按词元做 top-$K$ 选择。对于多种词元分布，MoD 以稠密模型 50\% 的浮点运算量达到相近的困惑度。

\question{3}{为什么光子计算可能从根本上改变大语言模型推理？目前有哪些局限？}
光子矩阵乘法利用光波干涉，以光速计算点积，每次乘法的能耗接近于零，因为光子不会使电阻发热。理论上有三方面好处。（1）\textbf{消除带宽墙}。光互连（optical interconnect）可在一根光纤上每秒传输数太比特，远超 HBM 的带宽限制。（2）\textbf{能效}。相比采用 CMOS 的 GPU，每次浮点运算的能耗约为其 $1/100$ 至 $1/1000$。（3）\textbf{时延}。光速传播消除了电子开关时延。目前的局限包括以下几项。（1）\textbf{数字信号到光信号的转换}。
% <<<PDFPAGE 52>>>
要把权重和激活值加载为光信号，仍需用到电子模数／数模转换器（ADC/DAC），转换慢且耗电。（2）\textbf{精度}。目前光子系统的等效精度为 4–6 位，足以用于推理，但需要量化。（3）\textbf{可编程性}。重新设置光学矩阵的元素比在 GPU 上加载新权重慢。（4）\textbf{可扩展性}。目前的光子芯片能实现 $64\times64$ 到 $512\times512$ 的小型矩阵乘法，远小于大语言模型的权重矩阵。据原书所述，截至 2025 年，Lightmatter 的 Passage 芯片是最先进的商用光子推理加速器。
\translatornote{Lightmatter 在 2025 年将 Passage M1000 定位为用于连接芯片的三维光子互连器件。参见 \href{https://lightmatter.co/press-release/lightmatter-unveils-passage-m1000-photonic-superchip-worlds-fastest-ai-interconnect/}{Lightmatter 的产品发布说明}。}

\question{4}{多词元预测如何提高训练和推理效率？}
多词元预测用 $k$ 个预测头来训练模型，各预测头比前一个多预测一步。第 1 个头预测 $t+1$，第 2 个头预测 $t+2$，以此类推。所有预测头共享主干网络的表征。

\textbf{训练收益}有三点。（1）每次前向传播产生更多梯度信号，因为模型同时优化 $k$ 个预测目标。（2）改善长程表征。要预测 $t+4$，模型必须捕捉超出局部共现的高层规律。（3）Meta 的实验表明，预测 3–4 个词元能提升基准测试表现和编程能力，较小的模型尤其受益。

\textbf{推理收益}在于，可利用新增的预测头进行自推测解码，无须进一步增大模型。用第 2 至第 $k$ 个头为 $t+2$ 到 $t+k$ 拟出草稿词元，然后用下一次主模型的前向传播验证。这样，无须另设草稿模型，也能获得 Medusa 式的加速。各头的接受率适中（约 0.7–0.8），在编程任务中可使解码加速 1.8–2.5 倍。

\question{5}{请说明什么是“推理优先”的模型设计理念。哪些架构选择体现了这种思路？}
推理优先（inference-first）从“生产部署会带来哪些约束？”这个问题出发，由部署约束反推架构决策。另一种做法是在训练时尽量提高基准测试分数，然后再设法优化推理。体现推理优先的架构选择包括以下方面。（1）\textbf{GQA、MQA 和 MLA}。减少 KV 头数可以显著降低 KV 缓存的显存占用，从而直接支持更大的批次和更长的上下文。（2）\textbf{采用 SwiGLU 的前馈网络}。按原书的描述，这种门控架构可按路由分数跳过混合专家变体中的整个专家，从而实现条件计算（conditional computation）。（3）\textbf{激活稀疏性（ReLU 变体）}。利用稀疏的激活模式，可以跳过前馈网络的大量计算，例如 Mistral 所称用 ReLU 达到的 89\% 前馈网络稀疏度。（4）\textbf{缩小词表}。较小的词表意味着嵌入表更小，采样更快。（5）\textbf{字节级或字符级切分}。可消除词元切分的瓶颈，并改善多语种处理。（6）\textbf{以深度优先于宽度}。更窄、更深的模型可能更难训练，但跨 GPU 并行的效率更高。Phi-3 和 Gemma 体现了推理优先的思路。它们刻意控制模型规模，把单位参数带来的质量放在首位，绝对基准分数居于次要地位。
\end{interview}
